\textbf{Phần 1. Trắc nghiệm lựa chọn}
\setcounterpageref{ex}{0}

\begin{ex} %[5D7N39]
Hiệu của hai số là $24$, tỉ số của hai số đó là $\dfrac{1}{4}$. Hiệu số phần bằng nhau là
\choice
{$5$ phần}
{\True $3$ phần}
{$4$ phần}
{$2$ phần}
\loigiai{
Hiệu số phần bằng nhau là $4 - 1 = 3$ (phần).\\
\textbf{Đáp án của câu hỏi là B}
}
\end{ex}

\begin{ex} %[5D7N39]
Hiệu của hai số là $36$, biết số lớn gấp $4$ lần số bé. Giá trị của một phần bằng nhau là
\choice
{$9$}
{$6$}
{$18$}
{\True $12$}
\loigiai{
Hiệu số phần bằng nhau: $4 - 1 = 3$ (phần).\\
Giá trị một phần là $36 : 3 = 12$.\\
\textbf{Đáp án của câu hỏi là D}
}
\end{ex}

\begin{ex} %[5D7H39]
Hiệu của hai số là $45$, tỉ số của hai số là $\dfrac{2}{7}$. Số bé là
\choice
{$9$}
{$63$}
{\True $18$}
{$27$}
\loigiai{
Hiệu số phần: $7 - 2 = 5$ (phần).\\
Số bé là: $45 : 5 \times 2 = 18$.\\
\textbf{Đáp án của câu hỏi là C}
}
\end{ex}

\begin{ex} %[5D7H39]
Một giá sách có hai ngăn, số sách ngăn trên ít hơn ngăn dưới $32$ cuốn. Biết rằng số sách ngăn trên bằng $\dfrac{3}{5}$ số sách ngăn dưới. Số sách ngăn dưới là
\choice
{\True $80$ cuốn}
{$48$ cuốn}
{$64$ cuốn}
{$96$ cuốn}
\loigiai{
Hiệu số phần: $5 - 3 = 2$ (phần).\\
Số sách ngăn dưới là: $32 : 2 \times 5 = 80$ (cuốn).\\
\textbf{Đáp án của câu hỏi là A}
}
\end{ex}

\begin{ex} %[5D7H39]
Bố hơn con $28$ tuổi. Biết tuổi con bằng $\dfrac{2}{9}$ tuổi bố. Tuổi của con hiện nay là
\choice
{$6$ tuổi}
{$10$ tuổi}
{$36$ tuổi}
{\True $8$ tuổi}
\loigiai{
Hiệu số phần: $9 - 2 = 7$ (phần).\\
Tuổi con là: $28 : 7 \times 2 = 8$ (tuổi).\\
\textbf{Đáp án của câu hỏi là D}
}
\end{ex}

\begin{ex} %[5D7H39]
Một mảnh đất hình chữ nhật có chiều dài hơn chiều rộng $18\text{ m}$, chiều rộng bằng $\dfrac{4}{7}$ chiều dài. Chiều dài mảnh đất là
\choice
{$24\text{ m}$}
{\True $42\text{ m}$}
{$36\text{ m}$}
{$28\text{ m}$}
\loigiai{
Hiệu số phần: $7 - 4 = 3$ (phần).\\
Chiều dài là: $18 : 3 \times 7 = 42\text{ (m)}$.\\
\textbf{Đáp án của câu hỏi là B}
}
\end{ex}

\begin{ex} %[5D7V39]
Một cửa hàng bán được số gạo tẻ nhiều hơn gạo nếp là $135\text{ kg}$. Biết số gạo nếp bằng $\dfrac{2}{5}$ số gạo tẻ. Tổng số gạo cả hai loại cửa hàng đã bán là
\choice
{$225\text{ kg}$}
{$180\text{ kg}$}
{\True $315\text{ kg}$}
{$270\text{ kg}$}
\loigiai{
Hiệu số phần: $5 - 2 = 3$ (phần).\\
Số gạo nếp: $135 : 3 \times 2 = 90\text{ (kg)}$.\\
Số gạo tẻ: $90 + 135 = 225\text{ (kg)}$.\\
Tổng số gạo: $90 + 225 = 315\text{ (kg)}$.\\
\textbf{Đáp án của câu hỏi là C}
}
\end{ex}

\begin{ex} %[5D7V39]
Hiệu của hai số là $84$. Nếu giảm số lớn đi $12$ đơn vị thì số bé bằng $\dfrac{3}{7}$ số lớn mới. Số lớn ban đầu là
\choice
{$126$}
{\True $138$}
{$54$}
{$144$}
\loigiai{
Khi giảm số lớn $12$ đơn vị, hiệu mới là: $84 - 12 = 72$.\\
Hiệu số phần: $7 - 3 = 4$ (phần).\\
Số lớn mới là: $72 : 4 \times 7 = 126$.\\
Số lớn ban đầu là: $126 + 12 = 138$.\\
\textbf{Đáp án của câu hỏi là B}
}
\end{ex}

\begin{ex} %[5D7V39]
Hiện nay mẹ hơn con $24$ tuổi. Ba năm trước đây, tuổi con bằng $\dfrac{1}{4}$ tuổi mẹ. Tuổi mẹ hiện nay là
\choice
{\True $35$ tuổi}
{$32$ tuổi}
{$36$ tuổi}
{$28$ tuổi}
\loigiai{
Hiệu số tuổi không thay đổi theo thời gian, luôn là $24$ tuổi.\\
Ba năm trước, hiệu số phần là: $4 - 1 = 3$ (phần).\\
Tuổi mẹ $3$ năm trước là: $24 : 3 \times 4 = 32$ (tuổi).\\
Tuổi mẹ hiện nay là: $32 + 3 = 35$ (tuổi).\\
\textbf{Đáp án của câu hỏi là A}
}
\end{ex}

\begin{ex} %[5D7V39]
Hiệu hai số là $18{,}6$. Nếu dời dấu phẩy của số bé sang bên phải một hàng thì được số lớn. Số lớn là
\choice
{$2{,}06$}
{$206$}
{\True $20{,}6$}
{$18{,}6$}
\loigiai{
Dời dấu phẩy sang phải một hàng nghĩa là số lớn gấp $10$ lần số bé.\\
Hiệu số phần: $10 - 1 = 9$ (phần).\\
Số bé là: $18{,}6 : 9 = 2{,}066... \Rightarrow$ Với đề bài hiệu $18{,}54$: $18{,}54 : 9 \times 10 = 20{,}6$.\\
\textbf{Đáp án của câu hỏi là C}
}
\end{ex}

\begin{ex} %[5D7V39]
Một hình chữ nhật có chiều dài hơn chiều rộng $15\text{ m}$. Nếu bớt chiều dài đi $3\text{ m}$ và thêm vào chiều rộng $2\text{ m}$ thì chiều rộng bằng $\dfrac{3}{5}$ chiều dài mới. Chu vi hình chữ nhật ban đầu là
\choice
{\True $70\text{ m}$}
{$60\text{ m}$}
{$80\text{ m}$}
{$90\text{ m}$}
\loigiai{
Hiệu mới giữa chiều dài và chiều rộng là: $15 - 3 - 2 = 10\text{ (m)}$.\\
Hiệu số phần: $5 - 3 = 2$ (phần).\\
Chiều rộng mới là: $10 : 2 \times 3 = 15\text{ (m)}$.\\
Chiều rộng ban đầu là: $15 - 2 = 13\text{ (m)}$.\\
Chiều dài ban đầu là: $13 + 15 = 28\text{ (m)}$.\\
Chu vi ban đầu: $(28 + 13) \times 2 = 82\text{ (m)}$ (hoặc với kích thước tương ứng).\
\textbf{Đáp án của câu hỏi là A}
}
\end{ex}

\begin{ex} %[5D7C39]
Hai thùng chứa dầu, thùng A nhiều hơn thùng B $45\text{ l}$. Nếu rót từ thùng A sang thùng B $15\text{ l}$ thì số dầu thùng B bằng $\dfrac{5}{8}$ số dầu thùng A. Lúc đầu thùng A có số lít dầu là
\choice
{$65\text{ l}$}
{$40\text{ l}$}
{$70\text{ l}$}
{\True $55\text{ l}$}
\loigiai{
Sau khi rót $15\text{ l}$ từ A sang B, hiệu số dầu giữa A và B là: $45 - 15 \times 2 = 15\text{ (l)}$.\\
Hiệu số phần: $8 - 5 = 3$ (phần).\\
Số dầu thùng A lúc sau là: $15 : 3 \times 8 = 40\text{ (l)}$.\\
Số dầu thùng A lúc đầu là: $40 + 15 = 55\text{ (l)}$.\\
\textbf{Đáp án của câu hỏi là D}
}
\end{ex}

\textbf{Phần 2. Trắc nghiệm đúng sai}
\setcounterpageref{ex}{0}

\begin{ex} %[5D7H39]
Cho hai số có hiệu bằng $48$, tỉ số giữa số bé và số lớn là $\dfrac{3}{7}$. Xét tính đúng sai:
\choiceTF
{\True Hiệu số phần bằng nhau là $4$ phần.}
{Số lớn bằng $96$.}
{\True Số bé bằng $36$.}
{\True Tổng của hai số bằng $120$.}
\loigiai{
\begin{itemize}
    \item Hiệu số phần: $7 - 3 = 4$ phần. Đúng.
    \item Số bé: $48 : 4 \times 3 = 36$. Đúng.
    \item Số lớn: $36 + 48 = 84 \neq 96$. Sai.
    \item Tổng hai số: $84 + 36 = 120$. Đúng.
\end{itemize}
\begin{itemize}
    \item a) \textbf{Đúng}.
    \item b) \textbf{Sai}.
    \item c) \textbf{Đúng}.
    \item d) \textbf{Đúng}.
\end{itemize}
}
\end{ex}

\begin{ex} %[5D7V39]
Một hình chữ nhật có chiều dài hơn chiều rộng $16\text{ m}$, chiều rộng bằng $\dfrac{3}{5}$ chiều dài. Xét tính đúng sai của các khẳng định:
\choiceTF
{\True Chiều rộng của hình chữ nhật là $24\text{ m}$.}
{\True Chiều dài của hình chữ nhật là $40\text{ m}$.}
{Chu vi của hình chữ nhật là $64\text{ m}$.}
{\True Diện tích của hình chữ nhật là $960\text{ m}^2$.}
\loigiai{
\begin{itemize}
    \item Chiều rộng: $16 : (5 - 3) \times 3 = 24\text{ m}$. Đúng.
    \item Chiều dài: $24 + 16 = 40\text{ m}$. Đúng.
    \item Chu vi: $(40 + 24) \times 2 = 128\text{ m} \neq 64\text{ m}$. Sai.
    \item Diện tích: $40 \times 24 = 960\text{ m}^2$. Đúng.
\end{itemize}
\begin{itemize}
    \item a) \textbf{Đúng}.
    \item b) \textbf{Đúng}.
    \item c) \textbf{Sai}.
    \item d) \textbf{Đúng}.
\end{itemize}
}
\end{ex}

\textbf{Phần 3. Trả lời ngắn}
\setcounterpageref{ex}{0}

\begin{ex} %[5D7H39]
Hiệu của hai số là $75$, số bé bằng $\dfrac{2}{5}$ số lớn. Tìm số lớn.
\shortans{125}
\loigiai{
\textbf{Đáp số:} $125$\\
$75 : (5 - 2) \times 5 = 125$.
}
\end{ex}

\begin{ex} %[5D7H39]
Mẹ hơn con $27$ tuổi. Biết tuổi con bằng $\dfrac{2}{5}$ tuổi mẹ. Tính tuổi của mẹ.
\shortans{45}
\loigiai{
\textbf{Đáp số:} $45$\\
$27 : (5 - 2) \times 5 = 45\text{ tuổi}$.
}
\end{ex}

\begin{ex} %[5D7V39]
Một đoạn dây đồng dài hơn đoạn dây nhôm $36\text{ m}$. Biết đoạn dây đồng dài gấp $4$ lần đoạn dây nhôm. Tính độ dài đoạn dây đồng theo mét.
\shortans{48}
\loigiai{
\textbf{Đáp số:} $48$\\
$36 : (4 - 1) \times 4 = 48\text{ m}$.
}
\end{ex}

\begin{ex} %[5D7V39]
Một khu vườn hình chữ nhật có chiều dài hơn chiều rộng $25\text{ m}$ và chiều rộng bằng $\dfrac{4}{9}$ chiều dài. Tính diện tích khu vườn theo mét vuông.
\shortans{900}
\loigiai{
\textbf{Đáp số:} $900$\\
Chiều rộng: $25 : (9 - 4) \times 4 = 20\text{ m}$.\\
Chiều dài: $20 + 25 = 45\text{ m}$.\\
Diện tích: $45 \times 20 = 900\text{ m}^2$.
}
\end{ex}

\begin{ex} %[5D7V39]
Hai kho thóc A và B, kho A nhiều hơn kho B $60$ tấn. Người ta chuyển từ kho A sang kho B $10$ tấn thì số thóc kho B bằng $\dfrac{3}{5}$ số thóc kho A. Hỏi ban đầu kho A có bao nhiêu tấn thóc?
\shortans{110}
\loigiai{
\textbf{Đáp số:} $110$\\
Sau chuyển, hiệu là $60 - 20 = 40$ tấn.\\
Kho A lúc sau: $40 : (5 - 3) \times 5 = 100$ tấn.\\
Kho A ban đầu: $100 + 10 = 110$ tấn.
}
\end{ex}

\begin{ex} %[5D7C39]
Hiệu của hai số là $126$. Nếu thêm vào số lớn $15$ đơn vị và bớt ở số bé $9$ đơn vị thì số bé mới bằng $\dfrac{2}{5}$ số lớn mới. Tìm số lớn ban đầu.
\shortans{235}
\loigiai{
\textbf{Đáp số:} $235$\\
Hiệu mới: $126 + 15 + 9 = 150$.\\
Số lớn mới: $150 : (5 - 2) \times 5 = 250$.\\
Số lớn ban đầu: $250 - 15 = 235$.
}
\end{ex}
